
\section{Klausurstruktur}

\begin{mytable}
    \begin{tabular}{@{}c l l@{}}
        \hline
        \textbf{Aufg.} & \textbf{Thema} & \textbf{Kernkompetenz}\\
        \hline
        1 & Sprachunterscheidung & Chomsky-Hierarchie, Pumping-Lemma, Grammatiken, RegExp\\
        2 & Endliche Automaten & NEA/DEA, Potenzmengenkonstr., Abschluss-Konstruktionen\\
        3 & Kontextfreie Gramm.\ \& LR/Kellerautomaten & Eindeutigkeit, LR-Parsing, Strukturbaum, CYK, PDA\\
        4 & Probleme / Berechenbarkeit \& Reduktion & P, NP, $\le_p$, Reduktionen, Chomsky-Größen\\
        5 & Multiple Choice & Querschnitt über den ganzen Stoff\\
        \hline
    \end{tabular}
\end{mytable}

\subsection{Sprachunterscheidung}
Zu einer Sprache $L$ entscheiden und beweisen, in welche Klasse sie fällt
(regulär / nicht regulär, kontextfrei / nicht kontextfrei).
Meist 3-4 Teilaufgaben, die die vier Beweisrichtungen kombinieren.
\begin{bem}
    Fast immer Varianten von $\{a^n b^m c^k \mid \text{lineare (Un-)Gleichung über } n,m,k\}$
\end{bem}

\subsubsection{Vollständige Klassifikation (reg/nreg und cf/ncf) mit Beweis.}
% \txc{2017H, 2017N, 2019H, 2019N, 2021, 2023P, 2023Z, 2025}

Für jede Sprache beide Fragen beantworten:
nicht regulär $\Rightarrow$ Pumping-Lemma (regulär); kontextfrei $\Rightarrow$ CFG/PDA;
nicht kontextfrei $\Rightarrow$ Pumping-Lemma (kontextfrei).

\begin{itemize}
    \item Ist die folgende Sprache regulär/nicht regulär, kontextfrei/nicht kontextfrei und warum:
    \begin{itemize}
        \item $\{\ol{x} \in \{\un{a}, \un{b}\}^* \mid  \ul{x} = a^n b^n \text{ mit } n \in \nz\} \cup \{ \ol{x} \in \{\un{c}, \un{d}\}^* \mid |\ol{x}| \leq 42 \}$ \txc{(HK2017)}
        \item
        $\{ \ol{z}  \in \{\un{a}, \un{b}, \un{c}\}^* \mid Anzahl(\un{a}, \ol{z}) = Anzahl(\un{b}, \ol{z}) = Anzahl(\un{c}, \ol{z}) \}$ \txc{(NK2017)}
        \item $\{a^m b^n c^k \mid n,m,k \in \nz_0, m=k \}$ \txc{(HK2019)}
        \item $\{a^m b^n c^k \mid n,m,k \in \nz_0, m \geq n \geq k \}$ \txc{(HK2019)}
        \item $\{(ab)^n \mid n \in \nz_0 \}$ \txc{(HK2019)}
        \item $\{a^m b^n \mid m+n  k; m,n,k \in \nz_0 \}$ \txc{(NK2019)}
        \item $\{a^m b^n c^k \mid n,m,k \in \nz_0, m=n \}$ \txc{(NK2019)}
        \item $\{a^m b^n c^k \mid n,m,k \in \nz_0, m \leq n \leq k \}$ \txc{(NK2019)}
        \item
        $\{\ol{x} \in \{\un{a}, \un{b},\un{c} \}^* \mid Anzahl(\un{a},\ol{x})+Anzahl(\un{b},\ol{x}) \text{ ist gerade } \land Anzahl(\un{b},\ol{x}) +Anzahl(\un{c},\ol{x}) \text{ ist ungerade }  \}$ \txc{(PK2021)}
        \item $\{a^m b^n c^k \mid n,m,k \in \nz_0, n = m+k\}$ \txc{(PK2021)}
        \item $\{ \ol{x} \in \{\un{a}, ..., \un{z}\}^* \mid \exists n: |\ol{x}| = 42^n \}$ \txc{(PK2021)}
    \end{itemize}
    \item Ist die folgende Sprache regulär/nicht regulär, kontextfrei/nicht kontextfrei? Beweisen Sie Ihre
    Antwort
    \begin{itemize}
        \item $\{\un{a}^n, \un{b}^m, \un{c}^k \mid n \leq m \leq k\}$ \txc{(PK2021)}
        \item  $\{a^n b^m c^k \mid n,m,k \in \nz_0, n+m = k \}$ \txc{(PK2023)}
        \item   $\{a^n b^m c^k \mid n,m,k \in \nz_0, n \leq m \} \cap \{ a^n b^m c^k \mid n,m,k \in \nz_{0}, m \leq k \}$
        \txc{(PK2023)}
        \item $\{ a^n b^m \mid n,m,k \in \nz_0, n+m = k, k > 5 \}$ \txc{(HK2023)}
        \item $\{ a^m b^n c \mid n,m \in \nz_0, n \leq m\} \cap \{a^n b^m c^k \mid n,m,k \in \nz_0, n\leq m \}$ \txc{(HK2023)}
        \item   $\{a^n b^m c^k \mid n,m,k \in \nz_0, n = k \} \cap \{ a^m b^n c \mid n,m \in \nz_0 \}$ \txc{(PK2025)}
        \item $\{a^{2025} b^{2025} c^k \mid k \in \nz_0, k \leq 2025 \}$ \txc{(PK2025)}
    \end{itemize}
\end{itemize}

\subsubsection{Nicht-Regularität mit dem Pumping-Lemma zeigen.} % \txc{2017H, 2017N, 2023P, 2023Z, 2025}
Typische Sprachen: $\{a^n b^m c^k\mid n=m\}$, $\{a^{2n}b^{4n}\}$, $\{c^4a^nb^n\}$.
Gefragt ist explizite Anwendung, das Gegenspieler-Schema muss also vollständig sein.

\begin{itemize}
    \item Zeigen Sie durch explizite Verwendung des Pumpinglemmas, dass die folgen-
    de Sprache nicht regulär ist:
    \begin{itemize}
        \item $\{\un{a}^{2n}, \un{b}^{4n} \mid n \in \nz\}$ \txc{(HK2017)}
        \item $\{ \un{c}^4, \un{a}^n, \un{b}^n \}$ \txc{(NK2017)}
        \item $\{a^n b^m c^k \mid n,m,k \in \nz_0, n = k \}$  \txc{(PK2023)}
        \item $\{a^n b^m c^k \mid n,m,k \in \nz_0, n = m \}$  \txc{(NK2023)}
        \item $\{a^n b^k c^m \mid n,m,k \in \nz_0, n = k\}$ \txc{(PK2025)}
    \end{itemize}
\end{itemize}

\subsubsection{Kontextfreiheit durch Angabe einer kontextfreien Grammatik zeigen.} % \txc{2017N, 2021, 2023P, 2023Z,2025}
Für $L_1$ aus dem Pumping-Teil wird meist direkt danach die passende CFG verlangt.

\begin{itemize}
    \item Geben Sie eine explizite kontextfreie Grammatik für L1 (Teilaufgabe ...)) an.
    \begin{itemize}
        \item $\{a^n b^m c^k \mid n,m,k \in \nz_0, n = k \}$  \txc{(PK2023)}
        \item $\{a^n b^m c^k \mid n,m,k \in \nz_0, n = m \}$  \txc{(NK2023)}
        \item  $\{a^n b^k c^m \mid n,m,k \in \nz_0, n = k\}$ \txc{(PK2025)}
    \end{itemize}
\end{itemize}


\subsubsection{Regularität durch regulären Ausdruck (oder DEA/rechtslin.\ Grammatik) zeigen.}% \txc{2017H, 2017N}
z.\,B.\ $\{x\in\{a,b\}^*\mid \#_a(x)\text{ gerade}\}$ oder $\#_a+\#_b\ge 2$.
Konstruktives Vorgehen: RegExp konkret hinschreiben.

\begin{itemize}
    \item Zeigen Sie durch die Konstruktion eines regulären Ausdrucks, dass die folgende Sprache regulär ist
    \begin{itemize}
        \item $\{ \ol{x} \in \{\un{a}, \un{b}\}^* \mid Anzahl(\un{a}, \ol{x}) gerade \}$ \txc{(HK2017)}
        \item
        $\{ \ol{x}  \in \{\un{a}, \un{b}, \un{c}\}^* \mid Anzahl(\un{a}, \ol{x}) = Anzahl(\un{b}, \ol{x}) \geq 2 \}$
        \txc{(NK2017)}
    \end{itemize}
\end{itemize}

\subsubsection{Schnitt/Vereinigung zweier Sprachen aka ``erst den Schnitt berechnen''.} % \txc{2023P, 2023Z, 2025}
\begin{itemize}
    \item $\{\ol{x} \in \{\un{a}, \un{b}\}^* \mid  \ul{x} = a^n b^n \text{ mit } n \in \nz\} \cup \{ \ol{x} \in \{\un{c}, \un{d}\}^* \mid |\ol{x}| \leq 42 \}$ \txc{(HK2017)}
    \item  $\{a^n b^m c^k \mid n,m,k \in \nz_0, n \leq m \} \cap \{ a^n b^m c^k \mid n,m,k \in \nz_{0}, m \leq k \}$
    \txc{(PK2023)}
    \item $\{ a^m b^n c \mid n,m \in \nz_0, n \leq m\} \cap \{a^n b^m c^k \mid n,m,k \in \nz_0, n\leq m \}$ \txc{(HK2023)}
    \item   $\{a^n b^m c^k \mid n,m,k \in \nz_0, n = k \} \cap \{ a^m b^n c \mid n,m \in \nz_0 \}$ \txc{(PK2025)}
\end{itemize}


\subsubsection{Grammatik für Vereinigung/Konkatenation zweier Sprachen.}
% \txc{2017H (CFG, $\cup$), 2017N (rechtslin., $\cdot$)}
Allgemeine Konstruktion $G_\cup$ bzw.\ $G_\bullet$ aus zwei gegebenen Grammatiken angeben.

\begin{itemize}
    \item Seien $G$ und $G'$ rechtslineare Grammatiken. Geben Sie eine rechtslineare Grammetik $G_{\circ}$ an, die
    genau $L(G) \circ L(G')$ erzeugt. \txc{(NK2017)}
    \item Seien $G$ und $G'$ kontextfreie Grammatiken. Geben Sie eine kontextfreie Grammetik $G_{\cup}$ an, die
    genau $L(G) \cup L(G')$ erzeugt. \txc{(HK2017)}
\end{itemize}

\subsubsection{Sprachen mit festen Schranken / endliche Sprachen.}
% \txc{2017H ($|x|\le 42$), 2021 ($|x|=42n$), 2025 ($a^{2025}b^{2025}c^k,\,k\le 2025$)}
Trickfrage: feste Exponenten bzw.\ obere Schranke $\Rightarrow$ die Sprache ist
endlich bzw.\ regulär, hier nicht auf das Pumping-Lemma hereinfallen!

\begin{itemize}
    \item $\{\ol{x} \in \{\un{a}, \un{b}\}^* \mid  \ul{x} = a^n b^n \text{ mit } n \in \nz\} \cup \{ \ol{x} \in \{\un{c}, \un{d}\}^* \mid |\ol{x}| \leq 42 \}$ \txc{(HK2017)}
    \item $\{ \ol{x} \in \{\un{a}, ..., \un{z}\}^* \mid \exists n: |\ol{x}| = 42^n \}$ \txc{(PK2021)}
    \item $\{a^{2025} b^{2025} c^k \mid k \in \nz_0, k \leq 2025 \}$ \txc{(PK2025)}
\end{itemize}

\begin{warn} $\ $
\begin{itemize}
    \item \textbf{Endlich $\Rightarrow$ regulär.} $a^{2025}b^{2025}c^{k},\,k\le2025$ hat nur endlich viele Wörter
    $\Rightarrow$ regulär (und damit kontextfrei). Die festen $2025$ sind Konstanten, keine Variablen!
    \item Das Pumping-Lemma ist \textbf{nur notwendig}, nicht hinreichend. Es beweist \emph{Nicht}-Regularität,
    \textbf{niemals} Regularität.
    \item ``Regulär $\Rightarrow$ auch kontextfrei'' nicht vergessen mit anzugeben, wenn nach beidem gefragt ist.
    \item Beim Schnitt-Typ: Klasse \emph{nach} dem Schnitt bestimmen. $\{n=k\}\cap\{\dots\}$ erzwingt oft
    $n=m=k$, was typischerweise \emph{nicht} kontextfrei ist.
\end{itemize}
\end{warn}

\subsection{Endliche Automaten}

\subsubsection{NEA aus rechtslinearer Grammatik konstruieren (als Diagramm).}%\txc{2019H, 2019N, 2021, 2023P, 2023Z,2025}
Nichtterminale $\to$ Zustände, ein zusätzlicher Finalzustand $f$ für Regeln $S::=a$.

\begin{itemize}
    \item Geben Sie zu folgender rechtslinearer Grammatik einen äquivalenten NEA (als Diagramm) an:
    \begin{itemize}
        \item \txc{(PK2021)}
        \item \txc{(PK2023)}
        \item \txc{(NK2023)}
        \item \txc{(PK2025)}
    \end{itemize}
    \item Betrachten Sie die rechtslineare Grammatik G. Vervollständigen Sie gemäß der Konstruktion in der Vorlesung
    folgenden NEA, der genau die Sprache erkennen soll, die die Grammatik erzeugt.
    \begin{itemize}
        \item \txc{(HK2019)}
        \item \txc{(NK2019)}
    \end{itemize}
    \item Geben Sie einen
\end{itemize}

\subsubsection{Potenzmengenkonstruktion NEA $\to$ DEA.}%\txc{2017H, 2017N, 2019H, 2019N, 2021, 2023P, 2023Z, 2025}
Fast jede Klausur... Mal als Diagramm, mal als Übergangstabelle. Müllzustand $\emptyset$ mitnehmen.

\subsubsection{DEA direkt für eine gegebene Sprache angeben.} %\txc{2017H, 2017N}
z.\,B.\ ``$\#_a(x)$ ungerade'' oder ``$\#_a(x)\equiv 1 \pmod 4$'' -- Zähl-Automat modulo\,$k$.

\subsubsection{Komplement-Automat konstruieren.} % \txc{2019H, 2019N, 2021, 2023P}
Finalzustände tauschen nur beim vollständigen DEA korrekt.

\subsubsection{Kleene-Stern $L(A)^*$ bzw.\ $L(A)\cdot\{a\}$ konstruieren.} %\txc{2017H, 2017N, 2019N, 2023P}
$\varepsilon$-Kanten von Finalzuständen zurück zum (neuen) Start; neuer akzeptierender Startzustand.

\subsubsection{Konkatenation $L(A_1)\cdot L(A_2)$ / Vereinigung $L(A_1)\cup L(A_2)$.}% \txc{2019H, 2023Z (Konkat.);2025 (Vereinig.)}
Diagramm \emph{und} allgemeine Tupel-Konstruktion (disjunkte Zustandsmengen $\Phi_1\cap\Phi_2=\emptyset$).

\subsubsection{Äquivalenz Grammatik $\leftrightarrow$ Automat beweisen/widerlegen.}%\txc{2021, 2023P, 2023Z, 2025}
$L(G_1)=L(G_2)$? oder $L(A)\subsetneq L(G)$? Gegenbeispielwort finden oder Inklusion beider Richtungen.

\section{Kontextfreie Grammatiken \& LR/Kellerautomaten}

\section{Probleme / Berechenbarkeit \& Reduktion}

\section{Multiple-Choice}

\subsubsection{Sprachen \& Automaten}
\begin{mytable}
    \begin{tabular}{@{}p{0.60\textwidth} c p{0.30\textwidth}@{}}
        \hline
        \textbf{Aussage} & \textbf{} & \textbf{Begründung}\\
        \hline
        Jede endliche Sprache ist regulär. & W & endlich $\Rightarrow$ RegExp als Aufzählung.\\
        $\exists$ Paar kontextfreier $L_1,L_2$ mit $L_1\cup L_2$ kontextfrei. & W & CFL sogar \emph{immer} unter $\cup$ abgeschl.\\
        $\exists$ Paar kontextfreier $L_1,L_2$ mit $L_1\cap L_2$ kontextfrei. & W & ``es gibt ein Paar'' (z.\,B.\ regulär); \emph{nicht} für alle.\\
        Kontextfreie Sprachen sind gegen Schnitt abgeschlossen. & F & Gegenbsp.\ $a^nb^nc^n$.\\
        $\exists$ kontextsensitive $L$ mit $L^C$ nicht kontextsensitiv. & F & kontextsensitiv unter Kompl.\ abgeschl.\\
        Reg.\ Sprache unendlich $\iff \exists$ Wort $n_0\le|w|<2n_0$ ($n_0=\#$Zust.). & W & Standard-Charakterisierung.\\
        Reg.\ Sprache unendlich, sobald ein Wort der Länge $\ge n_0$ existiert. & W & Zustand wiederholt sich $\Rightarrow$ pumpbar.\\
        Reg.\ Grammatik mit $n$ Nichtterm.\ $\Rightarrow$ jeder DEA $\ge n$ Zustände. & F & minimaler DEA kann kleiner sein.\\
        LR-/DPDA-Sprachen sind gegen Komplement abgeschlossen. & W & DCFL unter Komplement abgeschl.\\
        DPDAs/PDAs gegen Schnitt abgeschlossen. & F & schon CFL nicht unter $\cap$.\\
        DPDAs gegen Schnitt mit DEA abgeschlossen. & W & DCFL $\cap$ regulär $=$ DCFL.\\
        $\{a^nb^mc^nd\mid m,n\}$ durch DPDA erkennbar. & W & deterministisch: $a$ push, $c$ pop, $d$ prüfen.\\
        Für kontextfreie Sprachen ist das Wortproblem \emph{nicht} entscheidbar. & F & CYK entscheidet es ($\mathcal O(n^3)$).\\
        $G=(\{S\},\{1,2\},\{S::=S1S\mid 2\},S)$ ist eindeutig. & F & z.\,B.\ ``$21212$'' hat zwei Bäume.\\
        Jede links-/rechtslineare Grammatik ist eindeutig. & F & reguläre Grammatiken können mehrdeutig sein.\\
        $G:\ S::=SS\mid(S)\mid\varepsilon$ ist nicht eindeutig. & W & $\varepsilon$/$SS$ erzeugen mehrere Bäume.\\
        \hline
    \end{tabular}
\end{mytable}

\subsubsection{Berechenbarkeit \& Komplexität}
\begin{mytable}
    \begin{tabular}{@{}p{0.60\textwidth} c p{0.30\textwidth}@{}}
        \hline
        \textbf{Aussage} & \textbf{} & \textbf{Begründung}\\
        \hline
        Probleme in NP sind entscheidbar. & W & $\text{NP}\subseteq$ entscheidbar.\\
        Ein unentscheidbares Problem kann nicht in NP sein. & W & Kontraposition von oben.\\
        Ein unentscheidbares Problem kann nicht NP-hart sein. & F & z.\,B.\ Halteproblem ist NP-hart.\\
        Rekursiv aufzählbare Sprachen werden von det.\ TM akzeptiert. & W & det.\ TM $=$ nichtdet.\ TM (r.\,a.).\\
        $\{(a,b,c)\mid \varphi_{a+b}(c)\!\downarrow\}$ ist formale Regelsprache (r.\,a.). & W & Halteproblem ist semi-entscheidbar.\\
        $\{(a,b,c)\mid \varphi_{a+b}(c)\!\downarrow\}$ ist entscheidbar. & F & Halteproblem unentscheidbar.\\
        $\{(a,b)\mid \varphi_a(b)\!\uparrow\}$ ist rekursiv aufzählbar. & F & Komplement des Halteproblems, nicht r.\,a.\\
        Es gibt abzählbar viele verschiedene Turingmaschinen. & W & jede TM endlich kodierbar.\\
        Es gibt überabzählbar viele verschiedene Turingmaschinen. & F & siehe oben.\\
        $\text{P}\subseteq\text{PSPACE}$. & W & Standardinklusion.\\
        $\text{P}\subset\text{NP}\cap\text{co-NP}$. & W\,(?) & $\subseteq$ gilt; echte Inklusion offen.\\
        Menge der entscheidbaren Probleme gegen Schnitt abgeschlossen. & W & zwei Entscheider ``und''-verknüpfen.\\
        Für beliebige entscheidbare $L_1,L_2$ ist ``$L_1\subseteq L_2$'' entscheidbar. & F & reduziert von Leerheit/Halteproblem.\\
        Für beliebige kontextfreie $L_1,L_2$ ist ``$L_1\subseteq L_2$'' entscheidbar. & F & CFL-Inklusion unentscheidbar.\\
        Entscheidbar, ob eine TM nach $\le t$ Schritten hält. & W & einfach $t$ Schritte simulieren.\\
        Wortproblem für kontextsensitive Sprachen ist entscheidbar. & W & LBA / Chomsky-1.\\
        Formale Regelsprachen abgeschl.\ gegen $\cap,{}^C,\cup,*$. & F & Komplement fehlt (r.\,a.\ nicht unter $^C$).\\
        \hline
    \end{tabular}
\end{mytable}

\subsubsection{Reduktion, NP-Vollständigkeit, Gödelisierung}
\begin{mytable}
    \begin{tabular}{@{}p{0.60\textwidth} c p{0.30\textwidth}@{}}
        \hline
        \textbf{Aussage} & \textbf{} & \textbf{Begründung}\\
        \hline
        $A$ NP-vollständig $\iff A\le_p\text{SAT}$ und $\text{SAT}\le_p A$. & W & $A\le_p$SAT$\Rightarrow A\in$NP; SAT$\le_p A\Rightarrow$ NP-hart.\\
        $A$ NP-vollständig $\iff \text{SAT}\le_p A$. & F & $A\in$NP fehlt (das wäre nur NP-hart).\\
        $A$ NP-hart $\iff A\le_p\text{SAT}$. & F & Richtung falsch (müsste SAT$\le_p A$).\\
        $A$ NP-hart $\iff A\le_p\text{ILP}\wedge\text{ILP}\le_p A$. & F & beschreibt NP-\emph{vollständig}, nicht NP-hart.\\
        Für jede reguläre Sprache $A$ gilt $A\le_p$ 3-SAT. & W & regulär $\subseteq$ P $\subseteq$ NP; 3-SAT NP-vollst.\\
        $A\in\text{Chomsky-2}\Rightarrow$ 2-SAT $\le_p A$. & F & Gegenbsp.\ $A=\emptyset$ (kein Ziel für Ja-Instanzen).\\
        Longest-Path (Pfad Länge $\ge k$, knotendisjunkt) ist NP-hart. & W & Hamiltonpfad als Spezialfall.\\
        Erreichbarkeit ($\exists$ Pfad $u$--$v$) ist NP-hart. & F & in P (BFS/DFS).\\
        Fair teilen (Werte auf $k$ Personen gleich) ist NP-vollständig. & W & Partition/Bin-Packing.\\
        Bier auf Gläser ohne Überlauf verteilen. & W & Bin-Packing $\Rightarrow$ NP-hart.\\
        Wunsch-Zimmernachbar (je 1 Wunsch) erfüllbar? & F & Matching $\Rightarrow$ in P.\\
        $\{1,\dots,5\}^*\to\Nz,\ (n_1,\dots,n_t)\mapsto\sum n_i 10^i$ ist Gödelisierung. & W\,(?) & Ziffern $<10$ $\Rightarrow$ injektiv; def.-abhängig (nicht surjektiv).\\
        $\N^*\to\N,\ (n_1,\dots,n_t)\mapsto\sum n_i 10^i$ ist Gödelisierung. & F & unbeschr.\ $n_i$ $\Rightarrow$ Überträge, \emph{nicht} injektiv.\\
        Eine Gödelisierung muss polynomiell berechenbar sein. & F & nur berechenbar \& injektiv nötig.\\
        Universelle TM mit 42 Zuständen und 42 Bandzeichen. & W & UTMs existieren, $42\times42$ reicht üppig.\\
        TM mit 42 Zuständen simulierbar durch TM mit wenigen Zuständen. & W\,(?) & Zustands-/Zeichen-Tradeoff; \emph{Skript prüfen}.\\
        Universelle TM mit \emph{einem} Zustand. & (?) & umstritten -- \emph{im Skript / Tutorium klären}.\\
        \hline
    \end{tabular}
\end{mytable}